\documentclass[11pt]{article}
\usepackage[a4paper,margin=18mm]{geometry}
\usepackage[no-math]{fontspec}
\setmainfont{Latin Modern Roman}
\usepackage{amsmath,amssymb,amsthm,xfrac,xspace,hyperref,graphicx,comment}
\excludecomment{DefTralics}
\providecommand{\C}{}
\providecommand{\bibliofont}{\small}
\newcommand{\ie}{i.e.,\xspace}
\newcommand{\eg}{e.g.,\xspace}
\newcommand{\etc}{etc.\xspace}
\newcommand{\cf}{cf.\xspace}
\newcommand{\resp}{resp.\xspace}
\newcommand{\smaller}{\small}
\newcommand{\Subsection}{\subsection}
\newcommand{\Subsubsection}{\subsubsection}
\newcommand{\skpt}{\smallskip}
\emergencystretch=3em
\providecommand{\G}{}
\input{author-preamble.tex}

\begin{document}
\begin{center}{\Large Stable item 1930}\end{center}
\paragraph{Source and attribution.}
Daniel Bertrand, Bas Edixhoven, \emph{Pink’s conjecture on unlikely intersections and families of semi-abelian varieties}, Journal de l’École polytechnique — Mathématiques 7 (2020), publisher version, DOI 10.5802/jep.126. \href{https://jep.centre-mersenne.org/articles/10.5802/jep.126/}{Publisher article}. Authors' text: \href{https://creativecommons.org/licenses/by/4.0/}{CC BY 4.0}.
\paragraph{Editorial problem boundary.}
Prove only the single complex CM-elliptic Ribet-section construction of Section 2.4 and its consecutive relative counterexample conclusion: for nonzero anti-Rosati endomorphism difference, the canonical section curve is contained in no proper closed subgroup scheme of the Poincaré extension family and has infinitely many fibrewise torsion points. The canonical trivialisation, subgroup argument, and complete independent generalised-Jacobian torsion proof and compatibility in Section 6 are retained. No sharp torsion-order-attainment or higher-dimensional claim is selected.
\paragraph{Difficulty (editorial): H3.}
Canonical antisymmetric Poincaré trivialization creates a section above a nonzero Rosati-skew endomorphism. Generalized-Jacobian divisor norms and Weil reciprocity establish torsion specialization, while non-torsion extension classes force the cyclic section multiples to be Zariski dense in the full semiabelian family.
\paragraph{Editorial transcription note.}
Original author language, mathematical content and ordering are unchanged. Publisher front matter is replaced by this independent layout wrapper. Original native statement, proof and macros remain editable; bibliography is recovered from the matching cached publisher HTML. No assistant-generated proof or mathematical repair is supplied.
\clearpage
\input{author-body.tex}
\begin{thebibliography}{99}
\bibitem{Bertrand11} D. Bertrand - “Special points and Poincaré bi-extensions” , 2011, with an appendix by Bas Edixhoven
\bibitem{Bertrand13a} D. Bertrand - “Extensions panachées autoduales” , J. K-Theory 11 (2013) no. 2, p. 393-411
\bibitem{BMPZ} D. Bertrand, D. Masser, A. Pillay \& U. Zannier - “Relative Manin-Mumford for semi-Abelian surfaces” , Proc. Edinburgh Math. Soc. (2) 59 (2016) no. 4, p. 837-875
\bibitem{Bertrand18} D. Bertrand \& H. Schmidt - “Unlikely intersections in semiabelian surfaces” , Algebra Number Theory 13 (2019) no. 6, p. 1455-1473
\bibitem{BLR} S. Bosch, W. Lütkebohmert \& M. Raynaud - Néron models , Ergeb. Math. Grenzgeb. (3) , vol. 21 , Springer-Verlag, Berlin, 1990
\bibitem{Breen} L. Breen - “Biextensions alternées” , Compositio Math. 63 (1987) no. 1, p. 99-122
\bibitem{Chambert-Loir99} A. Chambert-Loir - “Géométrie d’Arakelov et hauteurs canoniques sur des variétés semi-abéliennes” , Math. Ann. 314 (1999) no. 2, p. 381-401
\bibitem{HodgeII} P. Deligne - “Théorie de Hodge. II” , Publ. Math. Inst. Hautes Études Sci. 40 (1971), p. 5-57
\bibitem{HodgeIII} P. Deligne - “Théorie de Hodge. III” , Publ. Math. Inst. Hautes Études Sci. 44 (1974), p. 5-77
\bibitem{Deligne-VS} P. Deligne - “Variétés de Shimura: interprétation modulaire, et techniques de construction de modèles canoniques” , in Automorphic forms, representations and $L$-functions (Corvallis, Ore., 1977), Part 2 , Proc. Sympos. Pure Math. , vol. XXXIII , American Mathematical Society, Providence, RI, 1979, p. 247-289
\bibitem{Faltings-Chai} G. Faltings \& C.-L. Chai - Degeneration of abelian varieties , Ergeb. Math. Grenzgeb. (3) , vol. 22 , Springer-Verlag, Berlin, 1990
\bibitem{Ferrand} D. Ferrand - “Conducteur, descente et pincement” , Bull. Soc. math. France 131 (2003) no. 4, p. 553-585
\bibitem{Gao-Pisa} Z. Gao - “A special point problem of André-Pink-Zannier in the universal family of Abelian varieties” , Ann. Scuola Norm. Sup. Pisa Cl. Sci. (5) 17 (2017) no. 1, p. 231-266
\bibitem{Howe} S. Howe - “Higher genus counterexamples to relative Manin-Mumford” , ALGANT Master Thesis, Univ. Leiden, July 2012, available at https://www.universiteitleiden.nl/binaries/content/assets/science/mi/scripties/howemaster.pdf
\bibitem{Husemoller} D. Husemöller - Elliptic curves , Graduate Texts in Math. , vol. 111 , Springer-Verlag, New York, 2004
\bibitem{Jacquinot-Ribet} O. Jacquinot \& K. A. Ribet - “Deficient points on extensions of abelian varieties by $\mathbb{G}_\mathrm{m}$” , J. Number Theory 25 (1987) no. 2, p. 133-151
\bibitem{Klingler} B. Klingler - “Hodge loci and atypical intersections: conjectures” , 2017
\bibitem{KUY} B. Klingler, E. Ullmo \& A. Yafaev - “Bi-algebraic geometry and the André-Oort conjecture” , in Algebraic geometry (Salt Lake City, 2015) , Proc. Sympos. Pure Math. , vol. 97 , American Mathematical Society, Providence, RI, 2018, p. 319-359
\bibitem{Lopuhaa} M. A. Lopuhaä - “Pink’s conjecture on semi-abelian varieties” , MSc. thesis, Univ. Leiden, September 2014, available at https://www.universiteitleiden.nl/binaries/content/assets/science/mi/scripties/lopuhaamaster.pdf
\bibitem{Moonen-L1} B. Moonen - “Linearity properties of Shimura varieties. I” , J. Algebraic Geom. 7 (1998) no. 3, p. 539-567
\bibitem{Moret-Bailly} L. Moret-Bailly - Pinceaux de variétés abéliennes , Astérisque , vol. 129 , Société Mathématique de France, Paris, 1985
\bibitem{MumfordAV} D. Mumford - Abelian varieties , Tata Institute of Fundamental Research Studies in Math. , vol. 5 , Hindustan Book Agency, New Delhi, 2008
\bibitem{Pink_thesis} R. Pink - Arithmetical compactification of mixed Shimura varieties , Bonner Math. Schriften , vol. 209 , Universität Bonn, Mathematisches Institut, Bonn, 1990, Dissertation, Rheinische Friedrich-Wilhelms-Universität Bonn,1989, https://people.math.ethz.ch/\textasciitilde{}pink/dissertation.html
\bibitem{Pink1} R. Pink - “A combination of the conjectures of Mordell-Lang and André-Oort” , in Geometric methods in algebra and number theory , Progress in Math. , vol. 235 , Birkhäuser Boston, Boston, MA, 2005, p. 251-282
\bibitem{Pink2} R. Pink - “A common generalization of the conjectures of André-Oort, Manin-Mumford, and Mordell-Lang” , 2005, preprint, https://people.math.ethz.ch/\textasciitilde{}pink/ftp/AOMMML.pdf
\bibitem{Ribet} K. A. Ribet - “Cohomological realization of a family of $1$-motives” , J. Number Theory 25 (1987) no. 2, p. 152-161
\bibitem{Serre} J.-P. Serre - Groupes algébriques et corps de classes , Publications de l’Institut Mathématique de l’Université de Nancago , vol. 7 , Hermann, Paris, 1984
\bibitem{SGA3.II} - Schémas en groupes. II: Groupes de type multiplicatif, et structure des schémas en groupes généraux (SGA 3.II) , Lect. Notes in Math. , vol. 152 , Springer-Verlag, Berlin-New York, 1970, Séminaire de Géométrie Algébrique du Bois Marie 1962/64. Dirigé par M. Demazure et A. Grothendieck
\bibitem{Tsimerman} J. Tsimerman - “The André-Oort conjecture for $\cal A_g$” , Ann. of Math. (2) 187 (2018) no. 2, p. 379-390
\bibitem{Zannier} U. Zannier - Some problems of unlikely intersections in arithmetic and geometry , Annals of Math. Studies , vol. 181 , Princeton University Press, Princeton, NJ, 2012
\end{thebibliography}
\end{document}
